%%%%%%%%%%%%%%%%%%%%%%%%%%%%%%%%%%%%%%%%%%%%%%%%%%%%%%%
% % Lines starting with % are comments, which are ignored.
% % This is a handy way of indicating the date and version of
% % your document, to wit:
% %
% % LaTeX sample file
% % Modified March, 2002
% %
%%%%%%%%%%%%%%%%%%%%%%%%%%%%%%%%%%%%%%%%%%%%%%%%%%%%%%%
\DocumentMetadata{
  % PDF standards: ONE key whose value is a comma-separated list --
  % UA-2 (accessibility) plus A-4f (archival; the archival companion of
  % UA-2). Do NOT use separate pdfstandard=... keys: with ordinary
  % key-value parsing only the last one takes effect.
  pdfstandard={UA-2,A-4f},
  lang=en-US,
  tagging=on,
  tagging-setup={
    math/setup=mathml-SE
  }
}

\documentclass[12pt]{exam}
% worksheet-macros.tex -- shared package setup for worksheet-template.tex.
%
% NOTE ON PATHS: every template in this repository is stored one level up
% from where it is compiled (see README.md). All relative \input paths
% below are resolved against the folder the document is compiled in,
% which sits one directory above this file.

% --- Math ------------------------------------------------------------
% amsmath: equation environments and numbering
% amsfonts: \mathbb etc.  |  mathrsfs: \mathfrak  |  stmaryrd: extra symbols
\usepackage{amsmath}
\usepackage{amsfonts}
\usepackage{mathrsfs}
\usepackage{stmaryrd}

% --- Layout and glyphs ------------------------------------------------
% float: [H] float placement (stop floats where you put them)
% fontspec: Unicode/OpenType fonts (compile with XeLaTeX or LuaLaTeX)
% emoji: emoji glyphs in the text
\usepackage{float}
\usepackage{fontspec}
\usepackage{emoji}

% --- Typography --------------------------------------------------------
% Atkinson Hyperlegible as the document-wide sans serif (sfdefault):
% a high-legibility face that is easier to read for students with
% dyslexia or low vision.
\usepackage[sfdefault]{atkinson}

% --- Colors -------------------------------------------------------------
% dvipsnames: extended named color palette (the Accessible* colors in
%   lecture-macros.tex are built on it);
% table: \rowcolors / \cellcolor for shading worksheet tables.
% NOTE: xcolor may only be loaded ONCE per document -- a second
% \usepackage{xcolor} with different options is a fatal "Option clash"
% error, so ALL of its options are listed here.
\usepackage[dvipsnames,table]{xcolor}

% --- Multimedia ---------------------------------------------------------
% animate: embeds audio/video (and legacy Flash) media in the worksheet.
\usepackage{animate}

% --- Bibliography --------------------------------------------------------
% biblatex (biber backend): citations and a reference list for worksheets.
\usepackage[
  backend=biber,
  style=alphabetic,
  url=true,
  doi=true,
  isbn=true,
  backref=true
]{biblatex}

% --- Links ---------------------------------------------------------------
% hyperref makes every \ref/citation/URL clickable in the PDF.
% Link colors must satisfy WCAG 1.4.3 (>= 4.5:1 contrast on white):
%   blue = 8.6:1, green!40!black = 7.2:1.
% (The old cyan/magenta values measured 1.25:1 / 3.1:1 and failed.)
\usepackage{hyperref}
\hypersetup{
  colorlinks=true,
  linkcolor=blue,
  filecolor=blue,
  urlcolor=green!40!black,
}

% --- Shared course macros ------------------------------------------------
% Root macros.tex: math shorthands used across all course documents
% (\RR, \NN, \ZZ, \QQ, \grad, \curl, \tr, \der, \pder, ...).
\newcommand{\RR}{\mathbb{R}}
\newcommand{\NN}{\mathbb{N}}
\newcommand{\ZZ}{\mathbb{Z}}
\newcommand{\QQ}{\mathbb{Q}}
\newcommand{\vv}{\vec{v}}
\newcommand{\ww}{\vec{w}}
\newcommand{\xx}{\vec{x}}
\newcommand{\yy}{\vec{y}}
\newcommand{\zz}{\vec{z}}
\newcommand{\uu}{\vec{u}}
\newcommand{\basis}[1]{\vec{e_{#1}}}
\newcommand{\id}{\mathrm{id}}
\newcommand{\CC}{\mathbb{C}}
\renewcommand{\div}{\operatorname{div}}
\newcommand{\grad}{\operatorname{grad}}
\newcommand{\curl}{\operatorname{curl}}
\newcommand{\tr}{\operatorname{tr}}
\newcommand{\der}[2]{\dfrac{d #1}{d #2}}
\newcommand{\pder}[3][]{\dfrac{\partial^{#1} #2}{\partial #3}}


% % Title and author(s)
%%%%%%%%%%%%%%%%%%%%%%%%%%%%%%%%%%%%%%%%%%%%%%%%%%%%%%%
\title{Worksheet 10 - Product and Quotient Rule, Trig Derivatives}
\author{MATH 1271}
\date{10/13/2026}
%%%%%%%%%%%%%%%%%%%%%%%%%%%%%%%%%%%%%%%%%%%%%%%%%%%%%%%
% %
% % The next command allows your in import encapsulated
% % postscript files, .epsf or .eps files, which
% % contain vector graphic image data.
% %
%%%%%%%%%%%%%%%%%%%%%%%%%%%%%%%%%%%%%%%%%%%%%%%%%%%%%%%
% %   The first thanks indicates your affiliation
% %
% %  Just the name here.
% %
% % Your mailing address goes at the end.
% %
% % \thanks is also how you indicate grant support
% %
%%%%%%%%%%%%%%%%%%%%%%%%%%%%%%%%%%%%%%%%%%%%%%%%%%%%%%%

%%%%%%%%%%%%%%%%%%%%%%%%%%%%%%%%%%%%%%%%%%%%%%%%%%%%%%%
% $ The next command ensures the margins are 1 inch on
% % all sides.
% %
%%%%%%%%%%%%%%%%%%%%%%%%%%%%%%%%%%%%%%%%%%%%%%%%%%%%%%%
\begin{document}
\newpage
\maketitle

\begin{enumerate}

\item Use relevant derivative rules to differentiate each function below.

\begin{enumerate}

\item $h(z)=\sqrt{z}+\dfrac{1}{z^4}+5^z$ (Please use Leibniz notation.)

\vfill

\item $r(x)=\dfrac{3^x}{x^4+1}$

\vfill

\item $p(t)=(t^3-t^2+t-1)(t^2+3t+1)$

\vfill

\item $g(y)=\dfrac{y^2e^y}{2y+1}$

\vfill

\end{enumerate}

\newpage



\item Suppose that $f$ and $g$ are differentiable functions such that
$$f(4)=2, \qquad f'(4)=6, \qquad g(4)=5, \qquad g'(4)=-3$$
Using this information, compute the derivatives of each of the following functions:

\begin{enumerate}

\item $h(x)=f(x)g(x)$. What is $h'(4)$?

\vfill

\item $S(t)=\dfrac{f(t)}{g(t)}$. What is $S'(4)$?

\vfill

%\item $j(y)=\dfrac{g(y)}{f(y)}$. What is $j'(4)$?

%\vfill

\item $q(z)=z^2f(z)+z$. What is $q'(4)$?

\vfill

\item $k(x)=\dfrac{x}{f(x)}$. What is $k'(4)$?

\vfill

\end{enumerate}

\newpage



\item Suppose that a bakery sells cakes. Each cake is priced at $p=20$ dollars. At this price, the company sells $500$ cakes per month and demand is falling at a rate of $30$ cakes (sold per month) per dollar. 

The bakery is looking over its finances. In particular, it is looking at its \emph{monthly revenue}:

$$\text{Monthly Revenue}=\Big(\text{Price per cake}\Big)\times\Big(\text{Number of cakes sold per month}\Big)$$

The bakery would like to increase its monthly revenue. To do this, the bakery is considering one of the following options:

\begin{itemize}
	\item increase the price of a cake to $p=21$ dollars;
	\item keep the price of a cake at $p=20$ dollars;
	\item decrease the price of a cake to $p=19$ dollars.
\end{itemize}

In this problem, we'll use calculus to help us determine what the bakery should do. Let $f(p)$ be the number of cakes sold (per month) when the price of a cake is set to $p$ dollars. 

\begin{enumerate}

	\item How can the statement ``when each cake is priced at $p=20$ dollars, the bakery sells $500$ cakes per month" be rephrased in terms of the function $f(p)$ and/or its derivative, $f'(p)$? 
	
	\vfill
	
	\item How can the statement ``when each cake is priced at $p=20$ dollars, demand is falling at a rate of $30$ cakes per dollar" be rephrased in terms of the function $f(p)$ and/or its derivative, $f'(p)$?
	
	\vfill
	
	\item Let $R(p)$ be the monthly revenue (in dollars) that the bakery earns each month from the sale of cakes. What is a formula for $R(p)$ in terms of $p$ and $f(p)$?
	
	\vfill
	
	\newpage
	
	\item Use the Product Rule to find a formula for $R'(p)$ in terms of $p$, $f(p)$, and $f'(p)$.
	
	\vfill
	
	\item At what rate is the monthly revenue changing when the price of a cake is set to $p=20$?
	
	\vfill
	
	\item Based on your answer to the previous part, what should the company do? Should it raise the price of a cake? Keep it the same? Or decrease the price of a cake? What would lead to higher monthly revenue? Why?

\vfill

\end{enumerate}
\newpage


\item Calculate each of the following derivatives:

\begin{enumerate}

\item $f(x)=(5x^3+\sin(x))(4^x-2\cos(x))$

\vfill

\item $p(t)=\dfrac{\cos(t)}{t^6\cdot 6^t}$

\vfill

\item $g(r)=\dfrac{r\sec(r)}{5^r}$

\vfill

\end{enumerate}

\newpage


\item Use the derivative rules that have been covered so far to answer each of the following questions about slope.

\begin{enumerate}

\item Find the (exact) slope of the tangent line to the graph of $y=2x+\dfrac{\sin(x)}{2}$ at the point where $x=\frac{\pi}{6}$. 

\vfill

\item Find the (exact) slope of the tangent line to the graph of $y=x^2\cos(x)$ at the point where $x=\frac{\pi}{2}$.

\vfill

\newpage

\item Let $f(\theta)=\theta-2\sin(\theta)$. On the interval $[0,2\pi]$, at which value(s) of $\theta$ does the graph of $f(\theta)$ have a horizontal tangent line?

\vfill

\item Let 
$$g(x)=\tan(x)+x, \qquad 0\leq x\leq 2\pi.$$
At which values of $x$ (between $0$ and $2\pi$) do we have $g'(x)=3$?

\vfill

\end{enumerate}


\newpage


\item The function 
$$P(t)=24+8\sin(t)$$
represents a population of a particular kind of animal that lives on a small island. Here, $P$ is measured in hundreds and $t$ is measured in decades since January 1, 2010. 

	\begin{enumerate}
	
	\item What is the instantaneous rate of change of $P$ on January 1, 2030? (Round to the nearest $0.01$ and include units on your answer.)
	
	\vfill

	\item Write a sentence in everyday language that explains how the population is behaving at this point in time.
	
	\vfill
	
	\item Was the population increasing or decreasing in size on January 1, 2025? At what rate?
	
	\vfill
	
	\item (A warm up for future classes:) Between January 1, 2010 and January 1, 2050, on what date (month and year) do we expect the population to be at its peak (as large as possible)?
	
	\emph{Hint:} Graph $P(t)$ on the interval $[0,4]$. Notice that at the peak population, the tangent line is horizontal (slope 0). At what value of $t$ will this happen? How do you translate this value of $t$ to a month and year?
	
	\vfill
	
	\end{enumerate}



\end{enumerate}

% Most people don't use these, so they are "commented out"
% by starting the lines with a "%"
%\begin{keywords}
%   \LaTeX, typesetting
%\end{keywords}

%\begin{AMS}
%   50C60, 18C25
%\end{AMS}

%%%%%%%%%%%%%%%%%%%%%%
% % Here is the start of the Text
%%%%%%%%%%%%%%%%%%%%%%
\end{document}
